% \DocumentMetadata{pdfstandard=A-2u}
\documentclass[a4paper,12pt]{amsart}
\input epsf
\def\useimages{1}
\usepackage{natbib}
\usepackage{color}
\usepackage{float}
\usepackage{epsfig}
\usepackage{hyperref}
\usepackage{graphicx}
\usepackage{subcaption}
\usepackage{caption}
\captionsetup[subfigure]{labelfont=rm} % isn't being used so, either use it or lose it cause it throws errors
\usepackage{verbatim}
\usepackage{stackengine}
\usepackage{mathrsfs}
\usepackage{amsmath}
\captionsetup[subfigure]{subrefformat=simple,labelformat=simple}
\renewcommand\thesubfigure{(\alph{subfigure})}
\DeclareMathAlphabet\mathbold{OML}{cmm}{b}{it}
\textheight=8.75in
\textwidth=6in
\oddsidemargin=-0.1in
\evensidemargin=-0.1in
\vsize 8.2in
\numberwithin{equation}{section}
\newtheorem{theorem}{Theorem}[section]
\newtheorem{lemma}{Lemma}[section]
\newtheorem{remark}{Remark}[section]
\newtheorem{corollary}{Corollary}[section]
\newtheorem{definition}{Definition}[section]
\newtheorem{algorithm}{Algorithm}[section]
\newtheorem{assumption}{Assumption}[section]
\newtheorem{proposition}{Proposition}[section]
\newtheorem{example}{Example}[section]
\newcommand{\curl}{\operatorname{curl}}
\renewcommand{\div}{\operatorname{div}}
\newcommand{\trace}{\operatorname{tr}}
\newcommand{\sign}{\operatorname{sign}}
\newcommand{\Range}{\operatorname{Range}}
\newcommand{\diag}{\operatorname{diag}}
\renewcommand{\span}{\operatorname{Span}}
\def\b1{{\mathbf 1}}
\def\bv{{\mathbf v}}
\def\bu{{\mathbf u}}
\def\bw{{\mathbf w}}
\def\bb{{\mathbf b}}
\def\bg{{\mathbf g}}
\def\bh{{\mathbf h}}
\def\br{{\mathbf r}}
\def\bs{{\mathbf s}}
\def\bd{{\mathbf d}}
\def\be{{\mathbf e}}
\def\bp{{\mathbf p}}
\def\bq{{\mathbf q}}
\def\bx{{\mathbf x}}
\def\by{{\mathbf y}}
\def\bz{{\mathbf z}}
\def\bn{{\mathbf n}}
\def\bbf{{\mathbf f}}
\def\bB{{\mathbf B}}
\def\bM{{\mathbf M}}
\def\bV{{\mathbf V}}
\def\bU{{\mathbf U}}
\def\bY{{\mathbf Y}}
\def\bF{{\mathbf F}}
\def\bA{{\mathbf A}}
\def\bD{{\mathbf D}}
\def\bH{{\mathbf H}}
\def\bN{{\mathbf N}}
\def\bT{{\mathbf T}}
\def\bP{{\mathbf P}}
\def\bQ{{\mathbf Q}}
\def\bS{{\mathbf S}}
\def\bR{{\mathbf R}}
\def\bepsilon{{\boldsymbol \epsilon}}
\def\balpha{{\boldsymbol \alpha}}
\def\bbeta{{\boldsymbol \beta}}
\def\bzeta{{\boldsymbol \zeta}}
\def\bxi{{\boldsymbol \xi}}
\def\btheta{{\boldsymbol \theta}}
\def\bpi{{\boldsymbol \pi}}
\newcommand{\Q}{{\mathcal Q}}
\newcommand{\R}{{\mathcal R}}
\newcommand{\A}{{\mathcal A}}
\newcommand{\B}{{\mathcal B}}
\newcommand{\C}{{\mathcal C}}
\newcommand{\D}{{\mathcal D}}
\newcommand{\F}{{\mathcal F}}
\newcommand{\G}{{\mathcal G}}
\newcommand{\I}{{\mathcal I}}
\newcommand{\J}{{\mathcal J}}
\newcommand{\M}{{\mathcal M}}
\newcommand{\N}{{\mathcal N}}
\newcommand{\X}{{\mathcal X}}
\newcommand{\Y}{{\mathcal Y}}
\newcommand{\calY}{{\mathcal Y}}
\newcommand{\calS}{{\mathcal S}}
\renewcommand{\L}{{\mathcal L}}
\renewcommand{\P}{{\mathcal P}}
\newcommand{\V}{\text{\bf V}}
\newcommand{\K}{{\mathcal K}}
\newcommand{\T}{{\mathcal T}}
\newcommand{\E}{{\mathcal E}}
\renewcommand{\O}{{\mathcal O}}
\newcommand{\hatcalK}{\widehat{\mathcal K}}
\newcommand{\hatcalS}{\widehat{\mathcal S}}
\newcommand{\hatA}{\widehat{A}}
\newcommand{\norm}[1]{\left\|#1\right\|}
\def\XVec#1{{\mathbf #1}}
\def\XNorm#1{\left\| #1 \right\|}
\def\XIProd#1#2{\left\langle #1 ,~ #2 \right\rangle}
\def\XM{\mu}
\def\XQ{Q}
\def\Xq#1{\XVec{q}_{#1}}
\def\Xu{\XVec{u}}
\def\Xf{\XVec{f}}
\def\Xe{\XVec{e}}
\def\Xr{\XVec{r}}
\def\Xx{\XVec{x}}
\def\Xv{\XVec{v}}
\def\Xw{\XVec{w}}
\def\Xn{\XVec{n}}
\def\Xes{\Xe_s}
\def\Xec{\Xe_c}
\def\Xus{\Xu_s}
\def\Xuc{\Xu_c}
\def\Xvs{\Xv_s}
\def\Xvc{\Xv_c}
\def\bone{{\boldsymbol 1}}
\def\bphi{{\boldsymbol \varphi}}
\def\bpsi{{\boldsymbol \psi}}
\def\btheta{{\boldsymbol \theta}}
\def\bchi{{\boldsymbol \chi}}
\def\boldeta{{\boldsymbol \eta}}
\def\bsigma{{\boldsymbol \sigma}}
\def\btau{{\boldsymbol \tau }}
\def\bdelta{{\boldsymbol \delta}}
\def\bQ{{\boldsymbol {Q}}}
\def\bI{{\boldsymbol {I}}}
\def\Nedelec{N\'ed\'elec\ }
\def\bPi{{\boldsymbol \Pi}}
\def\bPhi{{\boldsymbol \Phi}}
\def\bP{{\mathbf P}}
\usepackage{enumitem}
\newcommand{\ott}[1]{\bar{#1}}
% *** I changed a bunch of stuff directly, so definitely diff this.  But I left some things for you to sort out and those are marked by ***
\title[White Noise Sampling on Graphs]{White Noise Sampling on Graphs}
\author{Diya Patel and Panayot S. Vassilevski}
\address{Fariborz Maseeh Department of Mathematics and Statistics,
Portland State University, 1855 SW Broadway, Portland, OR 97201, U.S.A.}
\email{dpatel@pdx.edu, panayot@pdx.edu}
\keywords{fractional order hierarchical decompositions, multigrid,
fractional order white noise sampling}
\subjclass{65F10, 65N20, 65N30}
\date{May 29, 2026--beginning; Today is \today}
\thanks{This material is based upon work supported by the U.S. Department
of Energy, Office of Science, Advanced Scientific Computing Research Program.
The author benefited also from the Research Training Group (RTG) activities
under NSF grant DMS-2136228.}

\begin{document}

\begin{abstract}
This project extends by analogy a known Monte Carlo sampling procedure for 
discretized PDEs (partial differential equations) now to graph Laplacian 
problems. We implement a white noise sampling framework on graphs that 
generates spatially correlated random permeability fields on graph edges, 
solves a Darcy flow problem on the resulting weighted graph, and extracts a 
quantity of interest (QoI) that characterizes flow through the network. The 
framework is demonstrated on the Rhesus Macaque Brain Connectivity Network, 
where the QoI measures total neural signal flux from peripheral sensory 
regions to central processing hubs under uncertain synaptic connection 
strengths. We show that the Monte Carlo estimator converges reliably and 
that the resulting QoI distribution is approximately normal, consistent with 
theoretical predictions.
\end{abstract}

\maketitle

%=============================================================================
\section{Problem formulation}
%=============================================================================

We are given an undirected connected graph $G$ with vertices
$V=\{1,2,\;\dots\;n\}$ and a set of edges $E = \{e = (i,j) \} \subset V \times V$. 
The graph is uniquely described by its adjacency matrix $A = (a_{ij})$ with 
nonzero entries $a_{i,j}= a_{j,i} = 1$ if the vertices $i$ and $j$ form an 
edge $e = (i,j)$.

Let $d_i$ be the degree of vertex $i$ and $D = \diag(d_i)$ be the diagonal 
degree matrix. We also need the graph Laplacian matrix $L = D -A$, which can 
be represented in terms of the vectors $\bd_e \in {\mathbb R}^n$ for each 
edge $e=(i,j)$ with entries $1$ and $-1$ at positions $i$ and $j$ 
respectively. Namely, we have $L = \sum\limits_e \bd_e\bd^T_e$. We will be 
also interested in the weighted graph Laplacian for given edge-based 
weights $k =(k_e)_{e \in E}$
\begin{equation}\label{weighted graph Laplacian}
\sum\limits_{e \in E} k_e \bd_e \bd^T_e.
\end{equation}

The graph Laplacian $L$ plays a central role in this project. For any vector 
$\bv \in \mathbb{R}^n$, the quadratic form of $L$ satisfies:
\begin{equation}\label{quadratic form}
\bv^T L \bv = \sum_{e=(i,j) \in E} (v_i - v_j)^2.
\end{equation}
This identity shows that $L$ measures the total squared variation of $\bv$ 
across all edges. When $L$ appears in a linear system, the solution is forced 
to vary smoothly across connected vertices -- a property that underpins the 
spatial correlation structure generated in our sampling procedure.

First, we compute the minimal positive eigenvalue $\lambda^+_{\min}$ of 
$D^{-1} L$. For a connected graph, we have
\begin{equation*}
\lambda^+_{\min} = \min\limits_{\bv \perp \bone = (1)} \frac{\bv^T L \bv}{ \bv^T D \bv}.
\end{equation*}
This Rayleigh quotient characterization reveals the meaning of 
$\lambda^+_{\min}$: it is the minimum ratio of variation to connectivity 
across all non-constant vectors on the graph. The denominator $\bv^T D \bv$ 
normalizes by degree, ensuring that high-degree vertices do not artificially 
dominate. Intuitively, $\lambda^+_{\min}$ is the graph's fundamental 
frequency - the slowest non-trivial mode of variation consistent with the 
graph's connectivity structure. A small $\lambda^+_{\min}$ indicates a 
bottleneck in the graph; a large value indicates robust connectivity 
throughout.

Then, we form the shifted graph Laplacian
$\L_\sigma  = L + \sigma^2\lambda^+_{\min} D$
for a given parameter $\sigma^2$ (e.g., $\sigma^2 = 0.25$). We note that 
$\L_\sigma$ is s.p.d.(symmetric positive definite).

Our sampling procedure is as follows:
\begin{itemize}
\item 
We define a graph white noise $\bw = (w_e)_{e \in E}$ where $w_e$ are 
$\N(0,1)$ random numbers associated with each edge $e$. 

\item Then we form the vector $\bbf = (f_i)_{i=1}^{n}$ where $f_i$ is defined for 
every graph vertex $i$ as follows
\begin{equation*}
    f_i = \sum\limits_{e=(i,j)} w_e
\end{equation*}

\item The sampling process takes a random edge-based vector $\bw  = (w_e)$, 
forms the vertex-based vector $\bbf$ and computes the solution $\bu$ of the 
system of equations
\begin{equation}\label{shifted graph Laplacian white noise}
\L_\sigma \bu =  \lambda^+_{\min} \bbf.
\end{equation}
\end{itemize}

The sample $\bu =(u_i)$ corresponding to a $\bw$ is used as a random 
edge-based coefficient $k = (k_e)$, where 
\begin{equation}\label{random coefficient}
k_e = \exp\left({\frac{u_i+u_j}{2}}\right),\; e = (i,j).
\end{equation}

In a Monte Carlo method, we will solve at every Monte Carlo step a weighted 
graph Laplacian problem with the random edge-based coefficient $k$
\begin{equation}\label{weighted graph Laplacian problem}
\sum\limits_e k_e \bd_e \bd^T_e \bp = 0,
\end{equation}
subject to some values of $\bp$ prescribed on a given portion of vertices 
$\Gamma_{in}$ of the graph. We then compute a respective quantity of interest
(QoI), $Q(\bp)$, which depend on the solution $\bp$. More specifically, the 
QoI, $Q(\bp)$, is typically a function of a vector of entries of $\bp$ 
corresponding to another portion of graph vertices $\Gamma_{out}$ that are 
of particular interest. That is, we are given two vertex sets of interest, 
an input one, $\Gamma_{in}$, and an output one, $\Gamma_{out}$. We are 
interested how the input data propagates into the quantity of interest 
through the output values of the solution $\bp$ on $\Gamma_{out}$.

%=============================================================================
\section{Connection to the PDE Setting}
%=============================================================================

This project extends by analogy a well-established procedure for sampling 
spatially correlated Gaussian fields in the continuous PDE setting. In that 
setting, one generates a sample by solving a fractional-order elliptic PDE 
of the form:
\begin{equation}
(-\Delta + \kappa^2)^{\alpha/2} u = \mathcal{W},
\end{equation}
where $\Delta$ is the continuous Laplace operator, $\kappa > 0$ controls 
the correlation length, $\alpha > 0$ controls the smoothness of the 
resulting field, and $\mathcal{W}$ is spatial white noise. The solution 
$u$ is a Gaussian random field whose covariance function is a Mat\'ern 
kernel, characterized by exponential decay with distance. ***Put a citation here.***

The graph analogue replaces each component of this procedure as follows. 
The continuous Laplacian $-\Delta$ is replaced by the graph Laplacian $L$. 
The shift $\kappa^2$ is replaced by $\sigma^2 \lambda^+_{\min}$, calibrated 
to the graph's own spectral scale. The continuous spatial white noise 
$\mathcal{W}$ is replaced by graph white noise $w_e \sim \N(0,1)$ assigned 
independently to each edge. The fractional power $\alpha/2$ is set to $1$ 
in the current implementation, with extensions to fractional powers 
constituting future work.

The key insight is that solving $\L_\sigma \bu = \sqrt{\lambda_{\min}} \, \bbf$ 
on a graph produces a field with the same qualitative properties as the 
PDE solution: Gaussian marginals, spatial correlation that decays with 
graph distance, and a covariance structure tied to the graph's spectral 
properties. Specifically, since $\bbf = B\bw$ is a linear transformation 
of the Gaussian vector $\bw$, it is itself Gaussian. The solution 
$\bu = \L_\sigma^{-1} \lambda^+_{\min} \bbf$ is then also Gaussian, 
with covariance matrix:
\begin{equation}
\text{Cov}(\bu) = (\lambda^+_{\min})^2 \L_\sigma^{-1} B B^T \L_\sigma^{-T}.
\end{equation}
The covariance between two vertices $i$ and $j$ in $\bu$ depends on their 
graph-theoretic distance, decaying as vertices become more separated in the 
graph — the discrete analogue of the Mat\'ern correlation structure.

%=============================================================================
\section{Needed tools}
%=============================================================================

The project needs efficient graph Laplacian solvers, one for the shifted 
graph Laplacian problem \eqref{shifted graph Laplacian white noise} with 
white noise r.h.s. and another one, \eqref{weighted graph Laplacian problem}, 
for the weighted graph Laplacian with the computed random coefficient 
\eqref{random coefficient}. The two solvers have fundamentally different 
computational characteristics. Since $\L_\sigma$ is fixed across all Monte 
Carlo steps, its Cholesky factorization can be computed once and reused at 
every step. In contrast, $L_k$ changes at every step because $k$ changes, 
requiring a fresh solve each time. This asymmetry motivates different solver 
strategies for the two systems, as described in Section~\ref{sec:impl}.

%=============================================================================
\section{Implementation}\label{sec:impl}
%=============================================================================


\subsection{Phase 0: Data Import and Visualization}

Phase 0 essentially consists of two crucial parts: importing and 
visualizing network data.

\subsubsection{Importing the Data}
In this project, we are primarily dealing with unweighted and 
undirected graphs.

Data is imported by calling the \texttt{load\_graph()} function, which 
takes a filepath as input and returns an edge list, node count, and a 
dictionary of edge weights if present. The function handles two most 
common sources of network data files: Matrix Market (\texttt{.mtx}) 
files, parsed directly using \texttt{scipy.io.mmread}, and plain edge 
list text files, whose delimiter (whitespace or comma) and node 
indexing convention (0- or 1-indexed) is auto-detected from the file's 
content. Self-loops are discarded, and duplicate edges are removed (since we are dealing with undirected graphs). ***maybe rephrase as, "since we are treating all graphs as undirected at this stage"?***

Determining the total number of nodes is a crucial step here. One approach is to use the largest node index as an indication of the total number of node. However, this approach produces inaccurate results in many real-world datasets where non-contiguous node identifiers are used. For instance, during our research, we came across a graph that had node indices from 25 through 29, representing a 5-node graph. However, if we were to use the highest-node index approach here, we would end up falsely concluding that the graph has 30 nodes. In light of this, we used a relatively more robust approach to compute the total number of nodes in the graph in which we count the number of distinct nodes actually referenced across all edges to return the node count. Moreover, after collecting the full set of referenced node IDs, we re-map them to a contiguous range starting at 0, and report the true count of distinct nodes as n.

\texttt{load\_graph()} returns a dictionary consisting of edge weights 
in case of weighted graphs as one of its outputs even though this project specifically deals with unweighted graphs to avoid data loss. The inherent edge weights of a weighted graph are discarded by running 
\texttt{edge\_weights = None} before proceeding further with this 
pipeline, so that the edge permeabilities are derived entirely from the 
white noise sampling procedure described in Phases 3 and 4 of the pipeline and do not have any inherent bias to begin with.

\subsubsection{Visualizing the Data}

There is an optional but recommended step of visualizing the graph to 
get a visual sense of how the vertices and edges are laid out. The 
following four visualization functions can be used for this purpose:

\begin{itemize}

\item \texttt{visualize\_graph\_2d()} — generates a static 2D graph 
visualisation using Matplotlib and NetworkX. Node size and color reflect 
vertex degree. For weighted graphs, edge color reflects weight magnitude 
using the Plasma colormap. Node labels are shown automatically for graphs 
with up to 100 vertices. The figure can be saved to disk via 
\texttt{save\_as}, making it suitable for inclusion in the paper.

\item \texttt{visualize\_graph\_3d()} — generates a static 3D graph 
visualisation using Matplotlib. The $x$ and $y$ coordinates use the 2D spring 
layout; the $z$ coordinate uses the Fiedler vector (second eigenvector of 
the graph Laplacian), giving the third dimension structural meaning — 
vertices that are spectrally similar are placed at similar heights.

\item \texttt{visualize\_graph\_2d\_plotly()} — generates an interactive 
2D graph visualisation using Plotly. Fully zoomable and pannable, with 
hover tooltips showing vertex index and degree. Recommended for 
presentations and exploratory analysis.

\item \texttt{visualize\_graph\_3d\_plotly()} — generates a fully 
rotatable interactive 3D graph visualisation using Plotly. The $z$ 
coordinate uses the Fiedler vector for structural meaning. Most 
informative for presentations involving graphs with clear community 
structure. For very large graphs such as the US Power Grid 
($n = 4{,}941$), the 2D interactive version is preferred due to 
rendering cost.

\end{itemize}

It should be noted that for large graphs, visualization functions can 
be time-consuming. During our research, we also observed that based on the size of the graph, certain parameters that control the size of the nodes, color scale of the overall visuals, color of the edges, etc need to be modified to enhance visibility and interpretation of the visualisation. Also, certain network repositories, such as networkrepository.com, offer built-in visualization tools that can be used to get a quick visual sense of the data without having to build the visuals from scratch.

\subsection{Dataset: Oregon Router Network}

The Oregon Router dataset, sourced from the SNAP repository, represents 
a snapshot of Internet autonomous system (AS) peering relationships 
collected from Oregon route-views data. Each vertex corresponds to an 
autonomous system, and each edge represents a peering connection 
between two autonomous systems. After remapping to contiguous vertex 
IDs (Phase 0), the graph consists of $n = 11{,}174$ vertices and 
$23{,}409$ edges, with degrees ranging from 1 to 2{,}389 -- a wide 
spread, as shown in Figures~\ref{fig:oregon_network} and 
\ref{fig:oregon_degree_dist}.

\begin{figure}[htbp]
    \centering
    \includegraphics[width=0.85\textwidth]{oregon_network.png}
    \caption{Structure of the Oregon Router network ($n = 11{,}174$ 
    vertices, $23{,}409$ edges). Node size and color both reflect 
    vertex degree, which ranges from 1 to 2{,}389.}
    \label{fig:oregon_network}
\end{figure}

\begin{figure}[htbp]
    \centering
    \includegraphics[width=0.7\textwidth]{oregon_degree_dist.png}
    \caption{Degree distribution of the Oregon Router network (log 
    scale on the vertical axis). The distribution is heavily 
    right-skewed: most vertices have low degree, with a small number 
    reaching into the hundreds or above.}
    \label{fig:oregon_degree_dist}
\end{figure}


\subsection{Phase 1: Building the Graph Matrices}

The goal of Phase 1 is to map data into workable linear algebra format, 
where we form three matrices: adjacency matrix, degree matrix and 
graph Laplacian matrices. These matrices are denoted as $A$, $D$ and 
$L$ respectively going forward in this document.

The adjacency matrix, $A$, captures the relationship between edges and vertices 
of the graph and in case of unweighted graphs, has zeros and ones to 
indicate the absence or presence of edges between vertices and for 
weighted graphs, has numbers other than one to indicate the weights. 
For undirected graphs, $A$ is symmetric so, $a_{ij} = a_{ji}$ for all 
$i, j$. The degree matrix, $D$, is a diagonal matrix that captures the 
number of edges incident on every node, with diagonal entries 
$d_i = \sum_j a_{ij}$ for each $j$ (row sums of $A$). The 
Graph-Laplacian matrix, $L$, is one of the most crucial components of 
this research, which is computed by $L = D - A$.

The graph Laplacian is central to this project because for any vector 
$v$ belonging to $\mathbb{R}^n$, the quadratic form satisfies 
$v^T L v = \sum_{(i,j) \in E} (v_i - v_j)^2$, showing that $L$ measures the total squared variation of 
$v$ across all edges. A linear system involving $L$ produces solutions 
that vary smoothly across connected vertices, making $L$ the natural 
operator for enforcing spatial correlation in Phase 3. A key property 
that $L$ holds for a connected graph is that $L\mathbf{1} = 0$. Where $\mathbf{1}$ is a vector of all $1$'s of the appropriate length. This 
follows because $(D\mathbf{1})_i = d_i$ and 
$(A\mathbf{1})_i = \sum_j a_{ij} = d_i$ for every vertex $i$, so 
$D\mathbf{1} = A\mathbf{1}$, and therefore 
$L\mathbf{1} = (D - A)\mathbf{1} = 0$. This implies that $L$ 
has a zero eigenvalue, which makes it singular and non-invertible. 
This motivates modifying $L$ in Phase 2 and forming the shifted 
Laplacian which is symmetric positive definite, thus solving the 
singularity problem.

For relatively smaller graphs, $A$, $D$ and $L$ can be stored as dense 
NumPy arrays without running into any significant efficiency issues. 
However at scale, this becomes largely inefficient — an 
$11{,}174 \times 11{,}174$ dense matrix requires approximately 953 MB 
of RAM, whereas the same matrix when stored in sparse CSR format 
requires only a few megabytes since the vast majority of the entries 
are zeroes.

In code, this phase is implemented by calling the 
\texttt{build\_graph\_matrices()} function which takes the input of 
Phase 0's \texttt{load\_graph()} function (\texttt{edges}, 
\texttt{n\_vertices}). It starts by building $A$ first in a List of 
Lists (LIL) format which facilitates fast incremental construction, 
and then stores it in CSR format for the downstream arithmetic 
operations involving $A$. This flow has been designed such because LIL 
format is relatively cheaper for construction and CSR is better suited 
for the operations in the rest of the pipeline. $A$ is used to build 
$D$ by computing diagonal entries of $D$ as row sums of $A$ and having 
the off-diagonal entries as zeroes. $L$ is then computed using 
$D - A$. Then, all these matrices are efficiently stored as 
\texttt{scipy.sparse.csr\_matrix} objects.

The following sanity checks are performed on the output of this 
function to ensure that desired outputs are obtained:

\begin{enumerate}
    \item $A$ is symmetric: $A = A^T$
    \item $L$ is symmetric: $L = L^T$
    \item $L$ has a zero eigenvalue, confirmed by $L\mathbf{1} = 0$
    \item Dimensions of $A$, $D$ and $L$ are all $n \times n$ where 
    $n$ is the number of nodes.
\end{enumerate}

\subsection{Phase 2: The Shifted Laplacian}

Once the matrices are set up, we move on to phase 2. The goal of Phase 2 is to deal with the singularity of $L$. This is done by finding the smallest positive eigenvalue of $D^{-1}L$ and 
utilizing that to compute the shifted Laplacian $\L_\sigma$.

We need $L$ to be invertible so we can solve the system $Lu = f$ in 
Phase 3 to compute spatially correlated node noises. We make this 
possible by adding the smallest meaningful value to $L$ while also 
preserving its original connectivity structure. The shifted Laplacian 
is constructed using the relation,

\begin{equation}
\L_\sigma = L + \sigma^2 \lambda_{\min} D
\end{equation}
where $\L_\sigma$ is the shifted Laplacian — the invertible version 
of $L$. The components are:
\begin{itemize}[noitemsep, topsep=0pt]
    \item $L$ — the graph Laplacian from Phase 1
    \item $\sigma^2 = 0.25$ — a hyperparameter controlling the shift magnitude
    \item $\lambda^+_{\min}$ — the smallest positive eigenvalue of $D^{-1}L$
    \item $D$ — the degree matrix from Phase 1
\end{itemize}

Semantically, we do not want to add an arbitrary quantity to $L$ since that 
would fail to preserve the nature of $L$ and could distort the graph's 
connectivity structure. This is where $\lambda^+_{\min}$ comes into picture. 
It is the smallest most meaningful frequency of $L$ — the minimum ratio of 
variation to connectivity across all non-constant vectors on the graph, 
characterized by the Rayleigh quotient:
\begin{equation*}
\lambda^+_{\min} = \min_{\bv \perp \bone} \frac{\bv^T L \bv}{\bv^T D \bv}.
\end{equation*}
Using $\lambda^+_{\min}$ for the shift ensures the perturbation is calibrated 
to the graph's own spectral scale rather than being an arbitrary constant. 
We multiply by $D$ so that the shift is proportional to the degree of each 
vertex — higher degree vertices receive a larger bump than lower degree 
vertices, preserving the graph's geometric structure while eliminating the 
zero eigenvalue. The result $\L_\sigma$ is symmetric positive 
definite (SPD) — all eigenvalues are strictly positive — and therefore $\L_\sigma$ is
invertible.


\textbf{Implementation.} Phase 2 is implemented using two helper functions:

\begin{itemize}[noitemsep, topsep=4pt, leftmargin=3em]

\item \texttt{compute\_lambda\_min(L, D)} — computes $\lambda^+_{\min}$ 
by solving the generalized eigenvalue problem $L\bv = \lambda D\bv$ and 
returning the smallest positive result. The generalized form is preferred 
over computing $D^{-1}L$ directly because it preserves symmetry and is 
numerically stable. Two implementations are available: \texttt{scipy.linalg.eigh(L, D)}, 
a dense solver that computes all $n$ eigenvalues exactly — accurate but 
slow for large graphs since it computes eigenvalues we never use; and 
\texttt{scipy.sparse.linalg.eigsh(L, k=2, M=D, which='SM')}, a sparse 
iterative solver that finds only the 2 smallest eigenvalues. Since the 
graph is connected, the first eigenvalue is always zero and the second 
is $\lambda^+_{\min}$. This is significantly faster for large graphs at 
the cost of a small approximation error controlled by \texttt{tol=1e-10}, 
which is negligible in practice. For the US Power Grid ($n = 4{,}941$), 
\texttt{eigsh} reduces the eigenvalue computation from minutes to seconds. 
The function returns $\lambda^+_{\min}$ as a float.

\item \texttt{build\_shifted\_laplacian(L, D, lambda\_min, sigma\_squared=0.25)}
— takes $L$, $D$, and $\lambda^+_{\min}$ as inputs and computes 
$\L_\sigma = L + \sigma^2 \lambda^+_{\min} D$. A sanity check 
confirms that the smallest eigenvalue of $\L_\sigma$ is strictly 
positive, verifying that it is SPD and invertible. The function returns 
$\L_\sigma$ as output.

\end{itemize}

\begin{figure}[htbp]
    \centering
    \includegraphics[width=\textwidth]{phase2_eigenvalues.png}
    \caption{Eigenvalues of the generalized problem $(L,D)$ before 
    shifting (left) versus eigenvalues of the shifted Laplacian 
    $\L_\sigma$ after shifting (right), computed on a 15-vertex 
    illustrative graph. The unshifted problem has one eigenvalue at 
    exactly zero, corresponding to graph connectivity; this is what 
    makes $L$ singular. Adding $\sigma^2\lambda_{\min}D$ shifts every 
    eigenvalue upward by the same relative amount, eliminating the 
    zero eigenvalue and making $\L_\sigma$ symmetric positive definite.}
    \label{fig:phase2_eigenvalues}
\end{figure}


\subsection{Phase 3: White Noise Sampling}

This phase marks the beginning of constructing a spatially correlated 
edge permeability field for our flow network.

\subsubsection{3a. Draw Edge Noise}

In this step, draw as many random numbers from a standard normal 
distribution as there are edges in the graph. This can be denoted as vector $\mathbf{w}$ 
such that $\mathbf{w} \in \mathbb{R}^e$ where $e$ is the number of edges. A 
standard normal distribution is chosen because it produces zero-mean, 
unit-variance noise — ensuring the random field has no directional 
bias and that the magnitude of the noise is consistent regardless of 
the graph size. In code implementation, this is a single call to 
\texttt{numpy.random.randn(len(edges))}.

\subsubsection{3b. Map Noise to Nodes}

In this step, compute $f_i$ for each node $i$ by adding up $w_e$ of 
each edge $e$ incident on it. In matrix form, this is $\mathbf{f} = B\mathbf{w}$, where 
$B \in \mathbb{R}^{n \times e}$ for $n$ nodes and $e$ edges. This can 
be assumed to be a modified version of an incidence matrix where for 
each edge, its 2 nodes are assigned a value of 1. For larger graphs, 
$B \in \mathbb{R}^{n \times e}$ can be prohibitively large in dense 
format — for the Oregon Router Network ($n = 11{,}174$ and 
$e = 23{,}409$), a dense $B$ requires approximately 2.09 GB of RAM. In 
code, $B$ is built using the \texttt{build\_incidence\_matrix()} 
function, which constructs $B$ using LIL format and then converts it 
to CSR format for further computation. This is the same approach used 
to build $A$ in Phase 1. The result of this phase is a vector $\mathbf{f}$ such 
that $\mathbf{f} \in \mathbb{R}^n$ for $n$ nodes. The vector $\mathbf{f}$, however, has 
values that are uncorrelated. In practice, for the graphs and problems we are considering, neighboring nodes should have similar permeabilities. So, in the next 
step, we simulate the spatially correlated permeability field.

\subsubsection{3c. Make Noise Spatially Correlated}

In this step, we solve the system 
$\L_\sigma \mathbf{u} = \sqrt{\lambda_{\min}} \, \mathbf{f}$ for $\mathbf{u}$. This results in a 
vector $\mathbf{u}$ such that $\mathbf{u} \in \mathbb{R}^n$ for $n$ nodes. Since 
$\L_\sigma$ encodes graph connectivity, solving this system for $\mathbf{u}$ 
forces neighboring vertices to agree, thus producing a smooth $\mathbf{u}$ — a 
spatially correlated Gaussian random field where adjacent vertices 
have similar values, simulating the spatial correlation observed in 
real-world permeability fields.

Once $\mathbf{u}$ has been computed and the system 
$\L_\sigma \mathbf{u} = \sqrt{\lambda_{\min}} \, \mathbf{f}$ has to be solved, we would need 
to compute $\L_\sigma^{-1}$. However, at scale, computing the inverse 
of a matrix is computationally expensive and time consuming. To 
overcome this, we use Cholesky factorization. For a symmetric positive 
definite matrix such as $\L_\sigma$, Cholesky factorization decomposes 
it as $\L_\sigma = CC^T$, where $C$ is a lower triangular matrix. 
Solving $\L_\sigma \mathbf{u} = \sqrt{\lambda_{\min}} \, \mathbf{f}$ then reduces to two 
triangular substitutions — a forward solve $C\mathbf{y} = \mathbf{f}$ followed by a 
backward solve $C^T \mathbf{u} = \mathbf{y}$ — each step costing $O(n^2)$ as compared to 
the $O(n^3)$ cost of computing the matrix inverse.

The key advantage in the context of this project is that $\L_\sigma$ is 
fixed across all Monte Carlo sample runs. So, $\L_\sigma$ is factorized 
only once, and the factorized matrices are re-used in each run for two 
relatively cheaper solves.

Since Phase 3 has three sub-steps, a single function call to 
\texttt{run\_phase3()} runs the entire phase in one go. It takes the 
outputs of Phase 0 (\texttt{edges}, \texttt{n\_vertices}) and Phase 2 
(\texttt{L\_sigma} and \texttt{lambda\_min}) as input, and returns 
\texttt{u}, \texttt{f}, \texttt{B} and \texttt{factor}.

\textbf{Sanity Check.} To verify that Step 3c produces a smoother field 
than $\bbf$, we compare the difference between neighboring vertex values 
before and after the solve. For each edge $e = (i,j)$, we expect 
$|u_i - u_j| \ll |f_i - f_j|$. This is because $\bbf$ consists of sums 
of independent random numbers — neighboring vertices can have very 
different $f$ values since there is no constraint forcing them to agree. 
In contrast, $\bu$ is the solution of a linear system involving 
$\L_\sigma$, which penalizes large differences between connected 
vertices. As a result, $\bu$ exhibits significantly more uniform values 
across neighbors, directly reflecting the spatial correlation that 
$\L_\sigma$ enforces.



To visualize the spatial correlation introduced by the Phase 3c solve, 
we compare the value at each endpoint of every edge in the graph. 
Figure~\ref{fig:phase3_local_correlation} plots $f_i$ against $f_j$ 
(left) and $u_i$ against $u_j$ (right) for every edge in the Oregon 
Router network. If neighboring vertices' values were independent, no 
relationship would be visible in either scatter plot. Instead, the 
raw noise $f$ shows negligible correlation between edge endpoints 
(Pearson correlation $0.111$), while the solved field $u$ shows a 
strong, visually clear correlation ($0.616$) -- confirming that the 
Phase 3c solve produces genuine spatial correlation, with connected 
vertices tending to share similar values.

\begin{figure}[htbp]
    \centering
    \includegraphics[width=\textwidth]{phase3_local_correlation.png}
    \caption{Values at each edge's two endpoints, plotted against one 
    another, for the raw noise $f$ (left) and the spatially correlated 
    field $u$ (right), computed on the Oregon Router network 
    ($n=11{,}174$). Points clustering near the $y=x$ line indicate 
    that connected vertices tend to share similar values. The stark 
    contrast in correlation (0.111 versus 0.616) confirms that the 
    Phase 3c solve introduces genuine spatial correlation.}
    \label{fig:phase3_local_correlation}
\end{figure}


\subsection{Phase 4: Edge Permeability}

In Phase 4, we compute edge permeabilities for each edge based on the 
spatially correlated noise computed in Phase 3. For each edge, the 
following is computed,

\begin{equation}
k_e = \exp\!\left(\frac{u_i + u_j}{2}\right)
\end{equation}

The rationale behind this relation is that in a simple graph, each 
edge touches two nodes, and the $\exp$ function ensures that we get 
positive values only, since permeabilities cannot be negative. The 
result of this phase is a vector $\mathbf{k}$ such that 
$\mathbf{k} \in \mathbb{R}^e$ where $e$ is the number of edges.

\textbf{Implementation.} In code, it is implemented using a function called 
\texttt{compute\_permeability()} which takes two inputs: 
\texttt{u} (from Phase 3) and \texttt{edges} (from Phase 0). 
%Rather 
%than iterating over the edge list, the function extracts the endpoint 
%indices of all edges simultaneously into two arrays, and computes the 
%permeability for every edge in a single vectorized expression, 
%returning the result as an array \texttt{k\_vals}.

For efficiency, the function is fully vectorized using NumPy array indexing 
— all edge permeabilities are computed simultaneously in a single array 
operation rather than looping over edges one at a time:
\begin{equation*}
\mathtt{k\_vals} = \exp\!\left(\frac{\bu[\mathtt{i\_arr}] + 
\bu[\mathtt{j\_arr}]}{2}\right),
\end{equation*}
where \texttt{i\_arr} and \texttt{j\_arr} are precomputed integer arrays 
of the first and second endpoints of every edge respectively. This 
vectorization reduces Phase 4 computation time from $\mathcal{O}(|E|)$ 
Python iterations to a single NumPy call, which is particularly important 
for large graphs such as the Oregon Router Network ($|E| = 23{,}409$).

A sanity check confirms that all $k_e > 0$ — if this assertion fails it 
indicates a numerical issue upstream in Phase 3. After Phase 4, every edge 
in the graph has a permeability weight $k_e$, forming the weighted graph 
that is passed to the Darcy flow solve in Phase 5.


\subsection{Phase 5: Weighted Laplacian and Darcy Flow}

Phase 5 focuses on two important tasks: build the weighted Laplacian 
matrix ($L_k$) and use that to solve the Darcy flow problem to obtain 
the pressure field $\mathbf{p}$.

\subsubsection{5a. Build Weighted Laplacian}

In Phase 1, we built an unweighted graph Laplacian $L$ where each edge 
contributed an equal weight of 1. Later in Phase 4, we computed a 
random spatially correlated edge weight for each edge. To capture 
this, we build a weighted Laplacian, denoted as $L_k$, computed as:

\begin{equation}
L_k = \sum_{e = (i,j) \in E} k_e \, \mathbf{b}_e \mathbf{b}_e^T
\end{equation}

$L_k$ has a similar structure as $L$, but each edge contributes 
proportional to its edge permeability value ($k_e$). High permeability 
edges dominate the flow, whereas low permeability edges act as 
bottlenecks. Concretely, for an edge $e = (i,j)$ with 
permeability $k_e$, the incidence vector $\mathbf{b}_e$ produces 
a rank-one contribution to four entries of $L_k$:
\begin{equation}
L_k[i,i] \mathrel{+}= k_e, \quad
L_k[j,j] \mathrel{+}= k_e, \quad
L_k[i,j] \mathrel{-}= k_e, \quad
L_k[j,i] \mathrel{-}= k_e.
\end{equation}
Summing this contribution over all edges $e \in E$ yields the full 
weighted Laplacian $L_k = D_k - A_k$, where $A_k[i,j] = k_e$ is the 
weighted adjacency matrix and $D_k[i,i] = \sum_{e \ni i} k_e$ is the 
weighted degree. 

\textbf{Implementation.} $L_k$ is built using 
\texttt{build\_weighted\_laplacian(edges, n\_vertices, k)}, which takes 
the edge list and permeability dictionary from Phase 4 as inputs and 
returns $L_k$ as a sparse CSR matrix. For efficiency, $L_k$ is 
constructed in COO (Coordinate) sparse format rather than LIL format 
— COO format appends row indices, column indices, and data values to 
three arrays simultaneously, which is significantly faster than 
element-wise insertion for large edge sets.

Because every edge contributes to exactly four entries of $L_k$ (two 
diagonal, two off-diagonal, as shown above), we construct the COO 
representation by concatenating four aligned arrays — one entry per 
edge, per contribution:
\begin{verbatim}
row_idx = concat(i_arr, j_arr, i_arr, j_arr)
col_idx = concat(i_arr, j_arr, j_arr, i_arr)
\end{verbatim}
where \texttt{i\_arr} and \texttt{j\_arr} are the edge endpoint arrays. 
The first two blocks of \texttt{row\_idx}/\texttt{col\_idx} address the 
diagonal entries $(i,i)$ and $(j,j)$; the last two address the 
off-diagonal entries $(i,j)$ and $(j,i)$. The corresponding data array 
is built the same way, with $+k_e$ for the diagonal blocks and $-k_e$ 
for the off-diagonal blocks:
\begin{verbatim}
data = concat(k_vals, k_vals, -k_vals, -k_vals)
L_k = coo_matrix((data, (row_idx, col_idx)),
                  shape=(n_vertices, n_vertices)).tocsr()
\end{verbatim}
The COO constructor automatically sums contributions that map to the 
same $(row, col)$ position, which correctly accumulates the diagonal 
entries for vertices with multiple incident edges.

The sparsity structure of $L_k$ (\texttt{row\_idx} and \texttt{col\_idx}) 
is identical across all Monte Carlo steps since the graph topology does 
not change — only the permeability values $k_e$ change each step. This 
structure is therefore precomputed once before the Monte Carlo loop and 
reused at every step; only the \texttt{data} array (derived from the 
current sample's $k_e$ values) is rebuilt each iteration, reducing 
Phase 5a to a single array construction call per sample rather than a 
full re-derivation of the sparsity pattern.

A sanity check on $L_k$ verifies three properties:

\begin{enumerate}
    \item $L_k$ is symmetric
    \item All $k_e > 0$
    \item All row sums of $L_k$ are zero
\end{enumerate}

confirming that $L_k$ has the same structural properties as the 
unweighted Laplacian $L$.

\subsection{Gamma\_in / Gamma\_out Selection}
\label{sec:gamma_selection}

Choosing which nodes act as the source set $\Gamma_{\text{in}}$ and 
sink set $\Gamma_{\text{out}}$ is not governed by a single fixed rule — 
the appropriate method depends on the graph's underlying topology. 
Across the four datasets studied, three distinct selection methods were 
used, each chosen to match the structural characteristics of the 
corresponding graph rather than applying one approach uniformly.

\begin{table}[h]
\centering
\begin{tabular}{|l|l|l|}
\hline
\textbf{Approach} & \textbf{When it fits} & \textbf{Used for} \\
\hline
Degree-based & Clear hub-and-spoke structure & Rhesus Brain \\
\hline
Diameter endpoints & No dominant hub, real geometric extent & US Power Grid, Oregon Router \\
\hline
Community detection & Strong modular structure & bio-CE-GN \\
\hline
\end{tabular}
\caption{Gamma-selection method used per dataset, chosen to match each graph's dominant structural property.}
\label{tab:gamma_methods}
\end{table}

\subsubsection{Approach 1: Degree-Based Selection}

For graphs exhibiting a clear hub-and-spoke topology — a small number 
of high-degree hub vertices alongside a large population of low-degree 
peripheral vertices — the source and sink sets are defined directly 
from the degree distribution:
\begin{equation}
\Gamma_{\text{in}} = \{v : \deg(v) \leq D_{10}\}, \qquad
\Gamma_{\text{out}} = \{v : \deg(v) \geq D_{90}\},
\end{equation}
where $D_{10}$ and $D_{90}$ denote the 10th and 90th percentile of the 
degree distribution, respectively. This percentile-based cutoff was 
chosen over an arbitrary fixed count of vertices, since it provides a 
reproducible, self-documenting rule rather than an unmotivated choice 
of $k$. This method was applied to the Rhesus Brain dataset, whose 
degree distribution shows a clear separation between a small set of 
hub regions and many peripheral regions — directing flow from the 
periphery toward the hubs mirrors the graph's natural connective 
structure.

\subsubsection{Approach 2: Diameter Endpoints}

For graphs without a dominant hub structure but with meaningful 
geometric or topological extent — such as physical infrastructure 
networks — $\Gamma_{\text{in}}$ and $\Gamma_{\text{out}}$ are instead 
defined as neighborhoods around the two most topologically distant 
vertices in the graph. Computing the exact graph diameter is 
computationally prohibitive for large graphs, so the endpoints are 
approximated via a double breadth-first search: from a starting vertex, 
BFS ***version 1*** identifies the farthest vertex $v$; a second BFS from $v$ identifies 
the farthest vertex $w$. This is repeated from several starting 
vertices, retaining the pair with the greatest shortest-path distance. 
The source and sink sets are then taken as the $k$-hop neighborhoods 
around $v$ and $w$:
\begin{equation}
\Gamma_{\text{in}} = N_k(v), \qquad \Gamma_{\text{out}} = N_k(w),
\end{equation}
with any overlap between the two sets removed. This method was applied 
to the US Power Grid and Oregon Router datasets. We note that 
$|\Gamma_{\text{in}}|$ and $|\Gamma_{\text{out}}|$ can differ 
substantially even for identical $k$, since neighborhood size depends 
on local vertex density rather than distance alone; for Oregon Router, 
$k=4$ yielded $|\Gamma_{\text{in}}| = 11$ and $|\Gamma_{\text{out}}| = 92$. 
This asymmetry reflects genuine local structure in the graph and is 
corrected for at the comparison stage via normalization of $Q$ (see 
Section~\ref{sec:qoi_normalization}).

This approach towards gamma selection employs a double-sweep heuristic. The goal is to quantify flow in some sense from one point (could be a group of nodes for large graphs) to another so, we use this method to find two points that are farthest from each other in terms of number of hops. This heuristic proposes a breadth-first search (BFS)***version 2*** from an arbitrary starting vertex, which finds a fathest reachable vertex, and a second BFS from that vertex finds the true opposite diameter endpoint. This is a simplified variant of the double sweep lower-bound technique described by Crescenzi et 
al.~\citep{crescenzi2010finding}, who additionally propose a more 
sophisticated ``fringe'' algorithm to close the gap between this 
lower bound and the graph's true diameter -- a refinement not used in 
this work, where the double-sweep result alone is sufficient to 
identify two well-separated boundary regions.

\{Double-sweep heuristic.\} Approximating a graph's exact diameter requires computing the shortest paths between all pairs of nodes, which is obviously expensive for large graphs. ***version 3*** Instead, we use the double-sweep heuristic ~\citep{crescenzi2010finding} : from an arbitrary starting node $s$, a breadth-first search (BFS) finds the farthest reachable vertex $v$ ; a second BFS, starting from $v$, finds the farthest vertex from $v$, call it $w$. The distance between $v$ and $w$ is estimated to be the diameter of the graph.
    The first BFS serves only to locate a plausible extreme point of the graph -- since $s$ is an arbitrarily chosen node, it may lie anywhere relative to the graph's true diameter and so the distance between $s$ and $v$ would not be a reliable estimate for the graph's diameter.  This is where the second BFS comes in. It starts from $v$ and finds the node that is farthest from $v$, which we can call $w$. The distance between $v$ and $w$
can serve as a relatively more reliable estimate to be the graph's diameter since we are measuring distance between two nodes that we consider to be relatively extreme. In general $d(v,w) \geq d(s,v)$, since $v$ was itself the farthest 
point found from $s$; running the double sweep from several different 
arbitrary starting vertices and retaining the largest $d(v,w)$ found 
further improves the estimate.
\paragraph{Implementation.} \paragraph{Implementation.} In code, the double sweep is repeated from 
several different arbitrary starting vertices, keeping the pair 
achieving the largest distance $d(v,w)$ found across all attempts, 
since no single arbitrary starting vertex is guaranteed to yield the 
true diameter endpoints.
\subsubsection{Approach 3: Community Detection}

For graphs exhibiting strong modular structure — distinct, densely 
interconnected clusters rather than a hub-spoke or purely geometric 
topology — $\Gamma_{\text{in}}$ and $\Gamma_{\text{out}}$ are defined 
using community detection. Communities are identified via greedy 
modularity maximization \citep{clauset2004finding}, which begins with 
each vertex as its own community and iteratively merges the pair of 
communities yielding the greatest increase in modularity,
\begin{equation}
Q_{\text{mod}} = \frac{1}{2m}\sum_{ij}\left(A_{ij} - \frac{k_i k_j}{2m}\right)\delta(c_i, c_j),
\end{equation}
until no further merge improves the score. Clauset, Newman, and Moore 
note that community sizes found by this algorithm often follow a 
power-law distribution — a small number of large, dominant communities 
alongside many small ones — and observe this pattern directly in their 
own evaluation on a large-scale copurchasing network 
\citep{clauset2004finding}. We observe the same pattern across every 
dataset to which we applied community-based gamma selection: a small 
number of communities consistently account for the large majority of 
all vertices, with a long tail of much smaller communities. This is 
not an artifact specific to our datasets, but a documented, expected 
property of greedy modularity maximization, with direct consequences 
for the boundary fraction precondition discussed in 
Section~\ref{sec:boundary_fraction}.

The source and sink sets are then taken to be the two \emph{largest} 
detected communities:
\begin{equation}
\Gamma_{\text{in}} = C_{(1)}, \qquad \Gamma_{\text{out}} = C_{(2)},
\end{equation}
where $C_{(1)}$ and $C_{(2)}$ denote the largest and second-largest 
communities by vertex count. An earlier variant of this approach 
instead selected the community pair maximizing average inter-community 
shortest-path distance; this consistently favored small, structurally 
unrepresentative fragments over the graph's genuinely dominant modules; 
for example, an 18- and 15-vertex community pair was selected over the 
816- and 736-vertex dominant communities, since small peripheral 
fragments tend to be more distant in hop-count without reflecting 
meaningful structure. Selecting the two largest communities directly 
avoids this failure mode. This method was applied to the bio-CE-GN 
dataset, whose community structure is dominated by two large clusters 
(816 and 736 vertices) with a long tail of smaller communities — 
directly illustrating the power-law-like size distribution described 
above. ***after the dash sounds VERY AI to me***

We note that community-based selection can yield substantially larger 
$\Gamma_{\text{in}}$ and $\Gamma_{\text{out}}$ sets than the degree- or 
diameter-based approaches — hundreds of vertices, as opposed to a 
handful. Since $Q$ \ref{sec:qoi_normalization} is computed as a sum over 
edges incident to $\Gamma_{\text{out}}$, raw $Q$ values are not directly 
comparable across gamma-selection methods without normalizing by 
outlet edge count.

\subsubsection{Step 5b: Solve Darcy Flow}
\label{sec:darcy}

We then use the weighted graph Laplacian to solve the Darcy flow problem. 
The weighted Laplacian system encodes Darcy's law on the graph — for each 
edge $e = (i,j)$, the flow is $k_e|p_i - p_j|$, the discrete analogue 
of Darcy's law in porous media: flux equals permeability times pressure 
gradient. We prescribe boundary conditions on two disjoint sets of vertices:
\begin{itemize}[noitemsep, topsep=0pt, leftmargin=3em]
    \item $\Gamma_{in}$ — inlet vertices with prescribed pressure $p = 1$ 
    (signal source)
    \item $\Gamma_{out}$ — outlet vertices with prescribed pressure $p = 0$ 
    (signal sink)
\end{itemize}
The interior vertices $V \setminus (\Gamma_{in} \cup \Gamma_{out})$ are 
solved via the reduced system:
\begin{equation}
L_k \bp = \bb,
\end{equation}
where $\bb$ encodes the boundary conditions. The solution $\bp = (p_v)_{v \in V}$ 
is the pressure at every vertex. We verify that all interior pressure values 
satisfy $0 \leq p_v \leq 1$, confirming physically consistent flow behavior.

\textbf{Implementation.} The solver function used is:
\begin{itemize}[noitemsep, topsep=4pt, leftmargin=3em]
    \item \texttt{solve\_darcy(L\_k, n\_vertices, gamma\_in, gamma\_out)} 
    — uses\\ \texttt{scipy.sparse.linalg.spsolve}, a sparse direct solver. 
    Recommended for graphs up to approximately 10,000 nodes. The boundary 
    and interior vertex partitions are precomputed once before the Monte 
    Carlo loop and reused at every step to avoid redundant computation.
\end{itemize}
Since $L_k$ changes at every Monte Carlo step because $k$ changes, a 
fresh solve is performed each time — unlike $\L_\sigma$ which 
is prefactorized once in Phase 3. For the datasets tested in this project, 
\texttt{spsolve} with \texttt{use\_amg=False} was found to be faster than 
AMG due to the overhead of building the multigrid hierarchy at each step.


\subsection{Phase 6: Quantity of Interest}
\label{sec:qoi_normalization}

The goal of Phase 6 is to extract a single scalar quantity $Q$ that 
summarizes the total flow through the network for a given Monte Carlo 
sample.

As per Darcy's law, flow is proportional to the product of permeability 
and pressure difference. In our case, the flow along each edge 
$e = (i,j)$ is:
\begin{equation}
\text{flow}_{ij} = k_e \times |p_i - p_j|.
\end{equation}
The absolute value ensures the flow contribution is always non-negative 
regardless of edge orientation — without it, the pressure difference at 
outlet vertices would always be negative since $p_{\Gamma_{out}} = 0$ 
and neighboring interior vertices have $p > 0$, leading to a negative 
$Q$ which has no physical meaning.

The quantity of interest $Q$ is then the total flux exiting through 
$\Gamma_{out}$, computed by summing the flow contributions of all edges 
that are incident on at least one outlet vertex:
\begin{equation}
Q = \sum_{\substack{e=(i,j) \\ e \cap \Gamma_{out} \neq \emptyset}} 
k_e \, |p_i - p_j|.
\end{equation}
By conservation of flow, $Q$ equals the total flux entering through 
$\Gamma_{in}$ — what flows in must flow out. This can be verified as 
a sanity check by computing the inlet flux separately and confirming 
the two values are equal up to numerical tolerance.

For the Rhesus Macaque Brain Network, $Q$ measures the total neural 
signal flux reaching the central hub regions from the peripheral 
sensory regions under one realization of the random permeability field. 
Running this across $N$ Monte Carlo samples produces a distribution of 
$Q$ that characterizes uncertainty in signal propagation due to unknown 
synaptic connection strengths.

\textbf{Implementation.} Phase 6 is implemented in 
\texttt{extract\_qoi(p, k, edges, gamma\_out)}, which takes the pressure 
field $\bp$ from Phase 5, the permeability dictionary $k$ from Phase 4, 
the edge list, and the outlet vertex set as inputs, and returns $Q$ as 
a float. For efficiency, the function is fully vectorized using NumPy 
masking — an outlet mask identifies all edges incident on $\Gamma_{out}$ 
in a single array operation, and the flux sum is computed using 
\texttt{numpy.sum} rather than a Python loop over edges. This is 
particularly important for large graphs such as the Oregon Router Network 
where looping over 23,409 edges 1,000 times would add significant 
overhead.

\subsection{Phase 7: Monte Carlo Simulation}

This phase is all about running $N$ Monte Carlo simulations. We run 
Phases 3--6 in a loop to generate a random spatially correlated 
permeability field, solve for its Darcy flow, and extract $Q$ for 
each sample, then build a distribution to extract the mean and variance.

At each iteration $i = 1, \dots, N$:
\begin{enumerate}[noitemsep, topsep=4pt, leftmargin=3em]
    \item Draw fresh edge noise $\bw \sim \mathcal{N}(0,1)^{|E|}$ 
    (Phase 3a)
    \item Compute $\bbf = B\bw$ and solve $\L_\sigma \mathbf{u} = 
    \sqrt{\lambda^+_{\min}}\,\mathbf{f}$ for $\mathbf{u}$ (Phases 3b, 3c)
    \item Compute edge permeabilities $k_e = \exp((u_i + u_j)/2)$ 
    (Phase 4)
    \item Build $L_k$ and solve $L_k\bp = \bb$ for pressure field 
    $\bp$ (Phase 5)
    \item Extract $Q_i = \sum k_e|p_i - p_j|$ over outlet edges 
    (Phase 6)
\end{enumerate}
This produces an empirical distribution of $Q$ that characterizes 
uncertainty in flow through the network due to unknown connection 
strengths.

The expected value of $Q$ is estimated as:
\begin{equation}
\mathbb{E}[Q] \approx \frac{1}{N} \sum_{i=1}^{N} Q_i,
\end{equation}
with standard error $\sigma_Q / \sqrt{N}$. By the strong law of 
large numbers:
\begin{equation}
\frac{1}{N} \sum_{i=1}^{N} Q_i \xrightarrow{a.s.} \mathbb{E}[Q] 
\quad \text{as } N \to \infty,
\end{equation}
and the standard Monte Carlo error bound gives:
\begin{equation}
\left| \frac{1}{N} \sum_{i=1}^{N} Q_i - \mathbb{E}[Q] \right| = 
\mathcal{O}\!\left(\frac{\sigma_Q}{\sqrt{N}}\right),
\end{equation}
where $\sigma_Q$ is the standard deviation of $Q$. This 
$\mathcal{O}(N^{-1/2})$ convergence rate is independent of the 
problem dimension — a key advantage of Monte Carlo over deterministic 
quadrature methods, which suffer from the curse of dimensionality.

\textbf{Implementation.} The Monte Carlo loop is implemented in 
\texttt{monte\_carlo\_loop\_tqdm(L\_sigma, lambda\_min, edges, 
n\_vertices, gamma\_in, gamma\_out, edge\_weights, N, use\_amg)}, 
which takes $\L_\sigma$ and $\lambda^+_{\min}$ from Phase 2 
as inputs and returns an array $Q \in \mathbb{R}^N$ of QoI values. 
Several optimizations are applied to minimize per-sample computation 
time:
\begin{itemize}[noitemsep, topsep=4pt, leftmargin=3em]
    \item The Cholesky factorization of $\L_\sigma$, the 
    incidence matrix $B$, the edge index arrays, the COO sparsity 
    structure of $L_k$, and the boundary/interior vertex partition 
    are all precomputed once before the loop and reused at every step
    \item Phase 4 and Phase 6 are fully vectorized using NumPy array 
    operations, avoiding Python loops over edges
    \item A \texttt{use\_amg} flag allows switching between 
    \texttt{spsolve} (recommended for $n < 50{,}000$) and AMG 
    preconditioned CG (recommended for $n > 50{,}000$)
    \item A \texttt{tqdm} progress bar displays the current sample 
    count, elapsed time, estimated time remaining, and running mean 
    of $Q$
\end{itemize}

\textbf{Analysis.} The results are analysed using 
\texttt{analyze\_qoi(Q\_samples, save\_as)}, which prints summary 
statistics (mean, standard deviation, variance, min, max) and 
generates two plots: a histogram of the $Q$ distribution with the 
mean marked, and a convergence plot of the running mean of $Q$ versus 
number of samples. The convergence plot is the key diagnostic — a 
flat running mean confirms the estimator has converged and $N$ samples 
are sufficient for a reliable estimate.


\begin{figure}[htbp]
    \centering
    \includegraphics[width=\textwidth]{oregon_single_level_qoi.png}
    \caption{Distribution of $Q$ (left) and Monte Carlo convergence 
    (right) for single-level sampling on Oregon Router, $N=1000$ 
    samples. The running mean stabilizes well before $N=1000$.}
    \label{fig:oregon_single_level_qoi}
\end{figure}

\subsection{Two-Level Monte Carlo}
\label{sec:mlmc_integration}

While the Monte Carlo loop described above converges reliably, each 
sample requires a full solve on the fine graph, which can become 
computationally expensive for large $n$. Two-level Monte Carlo reduces 
this cost by introducing a smaller \emph{coarse graph} — built by 
merging groups of fine vertices into single coarse vertices — and 
combining a large number of cheap samples on the coarse graph with a 
smaller number of paired fine/coarse samples used to correct for the 
coarse approximation's bias. This section describes the aggregation 
procedure used to build the coarse graph, a diagnostic precondition 
that determines whether the method is viable for a given dataset, and 
the results of applying it across six datasets.

\subsubsection{Precondition: Boundary Fraction}
\label{sec:boundary_fraction}

Before attempting aggregation, we check what fraction of the graph is 
claimed by $\Gamma_{\text{in}} \cup \Gamma_{\text{out}}$:

\begin{equation}
\text{boundary fraction} = \frac{|\Gamma_{\text{in}}| + |\Gamma_{\text{out}}|}{n}
\end{equation}

This is a crucial diagnostic for whether two-level 
Monte Carlo is likely to succeed.  Aggregation will reduce the number of nodes in the 
%can only coarsen the 
non-boundary portion of the graph, regardless of the specific 
aggregation strategy used (Section~\ref{sec:aggregation}) -- if $\Gamma_{\text{in}}$ and 
$\Gamma_{\text{out}}$ together claim a large share of all vertices, 
the resulting coarse graph can end up with too few (or zero) interior 
vertices, making the reduced Darcy system in Phase 5b degenerate (no 
unknowns remain to solve for). We identify a 
concrete structural mechanism underlying this failure mode on 
hub-dominated graphs in Section~\ref{sec:rhesus}.

Across six datasets tested, we established the following empirical 
thresholds:

\begin{itemize}[noitemsep, topsep=4pt, leftmargin=3em]
    \item \textbf{Under 1\%:} aggregation reliably succeeds (US Power 
    Grid: 0.75\%, Oregon Router: 0.92\%, bio-grid-yeast with corrected 
    $k_{\text{hop}}$: 0.07\%)
    \item \textbf{10\% or more:} aggregation fails outright (Rhesus 
    Brain: 25\% $\to$ zero interior coarse vertices; bio-CE-GN: 70\%; 
    bio-grid-yeast via community detection: 76\%)
\end{itemize}

For diameter-endpoint $\Gamma$ selection specifically, boundary 
fraction depends on $k_{\text{hop}}$ \emph{relative to the graph's 
diameter}, not on $k_{\text{hop}}$ alone. A fixed $k_{\text{hop}}$ that 
works well on a large-diameter graph can fail catastrophically on a 
small-diameter, hub-dominated graph: $k_{\text{hop}}=4$ claimed only 
0.92\% of Oregon Router (diameter 10) but 99.87\% of bio-grid-yeast 
(diameter 5) — the same parameter, applied to a graph whose diameter 
was too small relative to it. Reducing $k_{\text{hop}}$ to 1 on 
bio-grid-yeast brought the boundary fraction down to 0.07\%, restoring 
a successful outcome.

This raises a natural question: how small should $k_{\text{hop}}$ be 
chosen? In general, $k_{\text{hop}}$ should be decreased until the 
boundary fraction falls comfortably below the empirical safety 
threshold established above. However, this tuning has a floor: at 
$k_{\text{hop}} = 0$, the neighborhood expansion is skipped entirely 
and $\Gamma_{\text{in}}$, $\Gamma_{\text{out}}$ reduce to the two 
diameter endpoint vertices themselves. If the boundary fraction 
remains too large even at this floor value, no further $k_{\text{hop}}$ 
tuning is possible, and the graph itself — rather than the choice of 
$k_{\text{hop}}$ — is the limiting factor.

We verified this directly on a small 15-vertex test graph: even at 
$k_{\text{hop}} = 0$ (using only the two diameter endpoint vertices as 
$\Gamma_{\text{in}}$ and $\Gamma_{\text{out}}$), the boundary fraction 
was $2/15 \approx 13.3\%$ — already an order of magnitude above the 
1\% threshold associated with reliable aggregation. This confirms that 
for sufficiently small graphs, two-level Monte Carlo via 
diameter-endpoint selection is not viable regardless of parameter 
tuning: the graph must be large enough that even a single-vertex 
boundary set on each side remains a small fraction of the total vertex 
count.

\subsubsection{Aggregation Algorithm}
\label{sec:aggregation}

Once boundary fraction has been checked and found acceptable, the 
coarse graph is built by grouping fine vertices into aggregates — one 
coarse vertex per aggregate — with the goal of substantially reducing 
vertex count while preserving enough of the fine graph's structure 
that the resulting coarse graph remains a useful approximation for 
paired sampling (Section~\ref{sec:paired_sampling}).

We considered several existing, documented graph coarsening 
strategies from the algebraic multigrid and graph partitioning 
literature, including independent-set-based approaches such as Luby's 
algorithm. These methods are well suited to coarsening a graph as a 
single, undifferentiated whole, but our setting imposes an additional 
structural requirement that general-purpose coarsening does not 
address: the source ($\Gamma_{\text{in}}$), sink ($\Gamma_{\text{out}}$), 
and interior vertex roles established during gamma selection must be 
preserved through coarsening, since the coarse graph inherits the same 
fixed-pressure boundary conditions used in the fine Darcy solve 
(Section~\ref{sec:darcy}). A coarse vertex that merged $\Gamma_{\text{in}}$ 
and $\Gamma_{\text{out}}$ vertices together, for instance, would be 
required to simultaneously hold pressure $1$ and pressure $0$ -- an 
ill-posed boundary condition. Consequently, the graph cannot be 
aggregated as a single undifferentiated unit; the aggregation 
procedure must be aware of, and respect, the source/sink/interior role 
of every vertex throughout. This requirement motivated the 
group-preserving algorithm described below, developed specifically for 
this constraint rather than adapted from an existing general-purpose 
method.

\textbf{Algorithm description.}Fine vertices are partitioned into 
three disjoint groups — $\Gamma_{\text{in}}$, $\Gamma_{\text{out}}$, 
and interior — and aggregated strictly within each group; a vertex is 
never paired with a neighbor from a different group. Within each 
group, vertices are processed in order of increasing degree. For each 
vertex, in turn: if it has an unassigned neighbor within the same 
group, the two are paired together as a new aggregate; otherwise, if 
it has a neighbor already belonging to an aggregate that has not 
exceeded a fixed size cap, the vertex joins that aggregate; otherwise, 
the vertex is left as a singleton aggregate. Processing vertices by 
increasing degree gives low-degree vertices, which have the fewest 
possible neighbors to pair with first opportunity to find a valid 
partner before their options are claimed by other vertices.

\begin{figure}[htbp]
    \centering
    \includegraphics[width=\textwidth]{oregon_fine_vs_coarse_bright.png}
    \caption{Fine graph (left) and coarse graph (right) for Oregon 
    Router, produced by the group-preserving aggregation algorithm. 
    Red vertices are $\Gamma_{\text{in}}$, cyan vertices are 
    $\Gamma_{\text{out}}$; gray (fine) and green (coarse) vertices are 
    interior. Node size in the coarse graph reflects aggregate size.}
    \label{fig:oregon_fine_coarse}
\end{figure}


Aggregating strictly within each group guarantees that $\Gamma_{\text{in}}$ 
and $\Gamma_{\text{out}}$ vertices can never merge into the same 
aggregate — a coarse vertex cannot simultaneously represent a fixed 
source ($p=1$) and a fixed sink ($p=0$) — without requiring this to be 
checked as a separate condition. The size cap on aggregate growth 
prevents a small number of aggregates from absorbing disproportionate 
edge permeability, which in earlier, uncapped variants of this 
algorithm inflated $Q_{\text{coarse}}$ by 10--100$\times$ relative to 
$Q_{\text{fine}}$ even after normalization.

\textbf{Coarse graph construction.} Regardless of aggregation 
strategy, every fine edge is checked against the aggregate assignment 
of its two endpoints. If both endpoints share an aggregate, the edge 
is internal and dropped; otherwise it contributes to a coarse edge 
between the two aggregates. A permanent mapping from each coarse edge 
to the fine edges that contributed to it is retained, since coarse 
edge permeability must be re-derived from that sample's fine 
permeability on every Monte Carlo sample, rather than sampled 
independently.

\subsubsection{Paired Sampling and the Two-Level Estimator}
\label{sec:paired_sampling}

For a single sample, white noise $w$ is drawn once and used to solve 
the fine white-noise system (Phase 3) in full, yielding $Q_{\text{fine}}$ 
via Phases 4--6. Coarse edge permeability is then derived, without 
drawing new randomness, by summing the fine permeability values 
belonging to each coarse edge according to the fixed mapping described 
above; the coarse Darcy system is solved identically to the fine one, 
yielding $Q_{\text{coarse}}$. Because both quantities originate from 
the same draw of $w$, they are correlated by construction — this 
correlation is what keeps $\mathrm{Var}(Q_{\text{fine}} - Q_{\text{coarse}})$ 
low enough that the correction term can be estimated accurately from a 
comparatively small number of expensive paired samples, while the bulk 
of the sampling budget is spent on relatively cheaper coarse-only samples for the 
first term of the estimator introduced earlier.

We note that coarse-only sampling is not as inexpensive as it may 
initially appear: since coarse permeability is derived from the fine 
field $u$, every coarse-only sample still requires the full fine 
white-noise solve (Phase 3) — only the fine Darcy solve (Phase 5b) is 
actually skipped. On this pipeline, where sparse Cholesky factorization 
has already made Phase 3 fast, this means the wall-clock savings from 
coarse-only sampling are more modest than the variance-reduction 
figures alone would suggest.

Two diagnostics determine whether the method is providing genuine 
benefit on a given dataset: the Pearson correlation between 
$Q_{\text{fine}}$ and $Q_{\text{coarse}}$ across paired samples, and 
the variance reduction ratio 
$\mathrm{Var}(Q_{\text{fine}})/\mathrm{Var}(Q_{\text{fine}}-Q_{\text{coarse}})$. 
Both are only meaningful once the underlying running means have 
converged; convergence behavior varies substantially by dataset (see 
Results), and was checked via running-mean plots rather than assumed 
from a fixed sample count in every case.

\subsubsection{A Mechanistic Failure Case: Rhesus Brain}
\label{sec:rhesus}

Beyond the boundary-fraction threshold itself, we identified a 
concrete structural mechanism by which a large boundary fraction 
causes aggregation to fail on hub-dominated graphs. On Rhesus Brain, 
using degree-based $\Gamma$ selection (top-decile degree as 
$\Gamma_{\text{in}}$, bottom-decile as $\Gamma_{\text{out}}$), we found 
that every one of the 68 interior vertices had \emph{zero} neighbors 
among other interior vertices — every edge touching an interior vertex 
connected exclusively to a boundary vertex. This renders the interior 
block $L_{ii}$ of the weighted Laplacian purely diagonal, decoupling 
the reduced Darcy system into 68 independent, single-variable equations 
with no diffusion between neighboring interior vertices. The resulting 
pressure field collapses to the source value ($p \approx 1$) 
uniformly across the interior (observed standard deviation 
$0.0000$), rather than showing the smooth gradient expected of a 
genuine flow solution. This confirms that a large boundary fraction on 
a hub-dominated graph does not merely reduce the interior vertex 
\emph{count} — a sufficiently connected hub can render the interior 
subgraph structurally disconnected from itself, independent of 
aggregation strategy.

\subsection{Results}

We evaluated the two-level Monte Carlo method across ***six/seven*** seven datasets 
spanning infrastructure, biological, and social network domains, 
using the gamma-selection strategy best suited to each dataset's 
structure (Section~\ref{sec:gamma_selection}).

\subsubsection{Successful Cases}
\hspace{1cm}        % makes the table show up after the subsubsection heading
\begin{table}[h!]
\centering
\begin{tabular}{lccccc}
\hline
Dataset & $n$ & Gamma method & Boundary frac. & Correlation & Var.\ reduction \\
\hline
US Power Grid    & 4{,}941  & Diameter-endpoint            & 0.75\% & 1.0000 & 3.11$\times$ \\
Oregon Router    & 11{,}174 & Diameter-endpoint            & 0.92\% & 0.9995 & 31.25$\times$ \\
bio-grid-yeast   & 6{,}008  & Diameter-endpoint ($k_{\text{hop}}{=}1$) & 0.07\% & 1.0000 & Strong \\
Facebook Ego     & 4{,}039  & Community (top-5\% subset)   & 2.20\% & 0.9956 & 12.66$\times$ \\
\hline
\end{tabular}
\caption{Two-level Monte Carlo results on datasets where the method succeeded.}
\label{tab:successes}
\end{table}

On bio-grid-yeast, the correction was necessary despite modest edge 
reduction (156{,}945 fine edges to 145{,}197 coarse edges, only 
7.5\% reduction) — a consequence of the graph's high density, which we 
discuss below. Despite this, correlation and variance reduction both 
remained strong, since the coarse solve's cost scales with vertex 
count, which was reduced 2.1$\times$, rather than edge count.

\subsubsection{Failure Cases}
\hspace{1cm}       % makes the table show up after the subsubsection heading
\begin{table}[h!]
\centering
\begin{tabular}{lccl}
\hline
Dataset & Boundary frac. & Correlation & Outcome \\
\hline
Rhesus Brain     & 25\% & --   & Zero interior coarse vertices \\
bio-CE-GN        & 70\% & --   & Predicted failure, not run \\
C.\ elegans      & 3.43\% & 0.72 & Structural checks pass; no statistical benefit \\
\hline
\end{tabular}
\caption{Datasets where two-level Monte Carlo failed, and why.}
\label{tab:failures}
\end{table}

\subsubsection{Structural Validity Is Necessary but Not Sufficient}

C.\ elegans is the most informative failure case, since it passes 
every structural precondition established above: boundary fraction 
(3.43\%) falls within the range associated with success, aggregation 
produces a substantial interior (92.5\% of coarse vertices), and no 
$\Gamma_{\text{in}}$/$\Gamma_{\text{out}}$ overlap occurs. Despite 
this correlation between $Q_{\text{fine}}$ and $Q_{\text{coarse}}$ was 
only $0.72$, and variance reduction was effectively $1.05\times$ — no 
meaningful benefit over direct fine-only sampling.

Comparing C.\ elegans against Facebook Ego Network — which used the 
identical gamma-selection strategy (a high-degree subset of the two 
largest detected communities) and has comparable graph density (mean 
degree $\approx 48$ vs.\ $\approx 44$) — isolates the difference to a 
single factor: the modularity of the underlying community structure.

\begin{table}[htbp]
\centering
\begin{tabular}{lccc}
\hline
Dataset & Modularity & Correlation & Var.\ reduction \\
\hline
C.\ elegans   & 0.26 (weak)   & 0.72  & 1.05$\times$ \\
Facebook Ego  & 0.78 (strong) & 0.996 & 12.66$\times$ \\
\hline
\end{tabular}
\caption{Modularity as a predictor of statistical success, holding gamma-selection method and graph density approximately constant.}
\label{tab:modularity}
\end{table}

We interpret this as follows: modularity measures how genuinely 
separated a detected community is from the rest of the graph. A 
low-modularity ``community,'' even when selected via the same greedy 
algorithm, remains substantially interconnected with the rest of the 
graph — the resulting $\Gamma_{\text{in}}$/$\Gamma_{\text{out}}$ do not 
correspond to a genuinely distinct flow regime, weakening the 
correspondence between the fine and coarse Darcy solutions even though 
every structural precondition is satisfied. We note that Clauset, 
Newman, and Moore's original description of greedy modularity 
maximization documents a power-law-like community size distribution — 
a small number of large, dominant communities alongside many smaller 
ones~\citep{clauset2004finding} - and we observe the same pattern 
directly: on bio-CE-GN and bio-grid-yeast, the two largest detected 
communities account for 70\% and 76\% of all vertices respectively, 
which is precisely why whole-community selection was replaced with a 
high-degree subset in later experiments.

This finding establishes that boundary fraction and interior 
non-degeneracy, while necessary, are not sufficient preconditions for 
two-level Monte Carlo to succeed when gamma selection is 
community-based; the modularity of the selected communities should be 
checked as an additional diagnostic.


\subsubsection{Comparison to Single-Level Monte Carlo}
\label{sec:single_vs_two_level}

To directly quantify the practical benefit of the two-level method, we 
compare it against single-level (fine-only) Monte Carlo on Oregon 
Router, matching both methods to approximately the same target 
standard error and reporting the number of expensive fine-graph 
samples required and the wall-clock time for each.

\begin{table}[htbp]
\centering
\begin{tabular}{lcccc}
\hline
Method & Fine samples & Total samples & Time (s) & SE \\
\hline
Single-level                    & 1000 & 1000 & 393.93 & 0.1092 \\
Two-level (this project)        & 100  & 2400 & 244.78 & 0.1072 \\
Two-level (general MLMC framework) & 100 & 2400 & 390.08 & 0.1086 \\
\hline
\end{tabular}
\caption{Single-level versus two-level Monte Carlo on Oregon Router, 
matched to approximately the same target standard error. The two-level 
row is reported both from this project's own implementation and from 
a general multi-level Monte Carlo framework developed alongside the 
PDE-based approach this project extends, via the adapter described in 
Section~\ref{sec:mlmc_integration}.}
\label{tab:single_vs_two_level_oregon}
\end{table}


\begin{figure}[htbp]
    \centering
    \includegraphics[width=\textwidth]{oregon_convergence.png}
    \caption{Convergence of $Q_{\text{fine}}$, $Q_{\text{coarse}}$ (left), 
    and the correction term $Q_{\text{fine}} - Q_{\text{coarse}}$ (right) 
    on Oregon Router, over 300 paired samples. Both quantities' running 
    means stabilize within approximately the first 100 samples, visually 
    confirming the high correlation ($0.9996$) and strong variance 
    reduction ($31.83\times$) reported above.}
    \label{fig:oregon_convergence}
\end{figure}


On Oregon Router, two-level Monte Carlo reached approximately the same 
precision (SE $=0.1072$ versus $0.1092$) using $10\times$ fewer 
expensive fine-graph samples (100 paired samples versus 1000 
single-level samples), and completed in $244.78\,$s versus 
$393.93\,$s for single-level -- a $1.6\times$ wall-clock speedup.

To validate this result independently, we re-ran the same comparison 
through a general multi-level Monte Carlo framework developed 
alongside the PDE-based approach this project extends 
(Section~\ref{sec:mlmc_integration}), using an adapter that implements 
the framework's required interface on top of this project's existing 
graph pipeline. At the same matched sample counts (2300 coarse-only, 
100 paired), the framework reached comparable precision (SE $=0.1086$) 
and a consistent estimate ($17.47$, versus $17.01$ from this project's 
own implementation and $17.5$ from the single-level benchmark), 
confirming the two implementations agree. The framework's run took 
$390.08\,$s, however -- noticeably longer than this project's own 
$244.78\,$s for the identical sample counts, and longer even than the 
single-level benchmark's $393.93\,$s. This overhead likely reflects 
the framework's additional per-sample bookkeeping (online variance 
tracking, structured result objects, and residual-norm verification 
inside its generic linear solver), which this project's own, more 
minimal implementation does not perform. The statistical benefit of 
two-level sampling -- fewer expensive samples needed for the same 
precision -- therefore transfers correctly across implementations, 
while the wall-clock benefit is implementation-dependent.

\begin{table}[htbp]
\centering
\begin{tabular}{lcccc}
\hline
Method & Fine samples & Total samples & Time (s) & SE \\
\hline
Single-level & 300 & 300  & 26.47  & 2.820 \\
Two-level    & 300 & 2300 & 119.47 & 2.608 \\
\hline
\end{tabular}
\caption{Single-level versus two-level Monte Carlo on US Power Grid, 
both using $N=300$ fine-graph samples for direct comparability. Unlike 
the Oregon Router comparison (Table~\ref{tab:single_vs_two_level_oregon}), 
sample counts here are not matched to a target precision; rather, both 
methods use the same fixed number of expensive fine-graph samples, and 
we compare the resulting precision directly.}
\label{tab:single_vs_two_level_powergrid}
\end{table}

On US Power Grid, at a fixed budget of 300 fine-graph samples, 
two-level Monte Carlo achieved a modestly better standard error 
($2.608$ versus $2.820$) than single-level sampling, but required 
substantially more wall-clock time ($119.47\,$s versus $26.47\,$s) -- 
the reverse of the trend observed on Oregon Router. This reflects the 
same underlying cause discussed above: coarse-only samples on this 
pipeline still pay the full Phase 3 cost, so the 2000 coarse-only 
samples used here ($\sim 93\,$s of the total) add substantial time 
without the corresponding precision benefit seen on Oregon Router, 
where the coarse Darcy solve's variance reduction was much larger 
(31.83$\times$ versus $3.59\times$ here). Combined with the heavy 
right-tailed behavior of $Q$ on this dataset (discussed above), 
US Power Grid illustrates that two-level Monte Carlo's benefit is not 
universal: its wall-clock advantage depends on both the achieved 
variance reduction and the dataset's convergence behavior, both of 
which varied substantially across the datasets tested in this work.

\subsection{Summary}

Across seven datasets, two-level Monte Carlo succeeded reliably when 
boundary fraction remained under approximately 1\% of total vertices, 
and failed when it exceeded roughly 25\%. For diameter-endpoint gamma 
selection, this threshold is governed by $k_{\text{hop}}$ relative to 
the graph's diameter rather than by $k_{\text{hop}}$ in isolation, and 
tuning $k_{\text{hop}}$ has a floor below which graph size itself 
becomes the limiting factor. For community-detection-based gamma 
selection, boundary fraction alone is insufficient to predict success: 
the modularity of the detected communities is an additional, necessary 
diagnostic, as demonstrated by the contrast between C.\ elegans and 
Facebook Ego Network. The Rhesus Brain case demonstrates a concrete 
mechanism by which large boundary fraction causes failure on 
hub-dominated graphs specifically: sufficiently high-degree boundary 
vertices can render the interior subgraph disconnected from itself, 
independent of the aggregation strategy used.

We additionally validated a second aggregation strategy -- 
group-preserving aggregation, which partitions vertices into 
$\Gamma_{\text{in}}$, $\Gamma_{\text{out}}$, and interior and 
aggregates strictly within each group -- against our original 
boundary-first strategy on US Power Grid and Oregon Router, finding 
the two strategies effectively interchangeable in terms of correlation 
and variance reduction. On Oregon Router, we further compared 
two-level Monte Carlo directly against single-level Monte Carlo at 
matched precision, finding a $10\times$ reduction in expensive 
fine-graph samples and a $1.6\times$ wall-clock speedup; the same 
comparison proved unreliable on US Power Grid and Facebook Ego Network 
due to the heavy right-tailed distribution of $Q$, itself a 
consequence of the log-normal edge permeability model used in Phase 4.

Several research questions raised by this work remain open. First, it 
is not yet established whether group-preserving aggregation's parity 
with boundary-first aggregation holds on datasets with large boundary 
fraction (e.g., Rhesus Brain) or merely on the well-behaved datasets 
tested so far. Second, the modularity threshold distinguishing 
successful from unsuccessful community-based gamma selection has only 
been characterized qualitatively (via a single contrasting pair of 
datasets) and not established quantitatively across a wider range of 
modularity values. Third, the heavy-tailed behavior of $Q$ complicates 
standard-error-based comparisons between sampling methods, and a more 
robust convergence and comparison methodology -- for instance, based 
on the log-transformed distribution, or on substantially larger 
benchmark runs -- remains to be developed. Fourth, the Minnesota Road 
Network's convergence instability was never fully diagnosed and 
remains an open case distinct from the boundary-fraction and 
modularity failure modes identified elsewhere in this work. Finally, 
integration with a collaborator's general multi-level Monte Carlo 
framework demonstrated that this pipeline's two-level structure is the 
$L=1$ special case of a fully general telescoping-sum estimator; 
whether the boundary-fraction and modularity findings established here 
extend to a genuine multi-level (three or more level) hierarchy, built 
by recursively coarsening the coarse graph, has not yet been 
investigated.
%=============================================================================
\section{Further extensions}
%=============================================================================

Transition to two- and multilevel Monte Carlo simulations 
(cf., \citet{Cliffe11}, \citet{Osborn17, Osborn17b}), and eventually 
to Markov Chain Monte Carlo (MCMC) and its multilevel extensions 
(cf., \citet{Dodwell15}, \citet{fairbanks}). Additional planned 
extensions include:
\begin{itemize}
\item \textbf{Fractional order sampling:} Extending to fractional 
powers $\alpha \neq 1$ to control the smoothness of the sampled 
field, more closely matching the continuous PDE setting.
\item \textbf{Sparse and multilevel solvers:} Replacing direct 
solvers with algebraic multigrid (AMG) preconditioned conjugate 
gradient for scalability to larger graphs beyond the current 
dense solver limits.
\end{itemize}

\bibliographystyle{plainnat}
\bibliography{betterrefs}
%=============================================================================
% \begin{thebibliography}{9}

% \bibitem{Osborn17}
% S.~Osborn, P.~Vassilevski, and U.~Villa,
% \textit{A multilevel, hierarchical sampling technique for spatially 
% correlated random fields},
% SIAM Journal on Scientific Computing, 39(5):S543--S562, 2017.

% \bibitem{clauset2004finding}
% A.~Clauset, M.E.J.~Newman, and C.~Moore,
% \textit{Finding community structure in very large networks},
% Physical Review E, 70(6):066111, 2004.

% \bibitem{Cliffe11}
% K.A.~Cliffe, M.B.~Giles, R.~Scheichl, and A.L.~Teckentrup,
% \textit{Multilevel Monte Carlo methods and applications to elliptic 
% PDEs with random coefficients},
% Computing and Visualization in Science, 14(1):3--15, 2011.

% \bibitem{Osborn17b}
% S.~Osborn, P.~Vassilevski, and U.~Villa,
% \textit{Smoothed aggregation for Helmholtz problems},
% Numerical Linear Algebra with Applications, 2017.

% \bibitem{Dodwell15}
% T.J.~Dodwell, C.~Ketelsen, R.~Scheichl, and A.L.~Teckentrup,
% \textit{A hierarchical multilevel Markov chain Monte Carlo algorithm 
% with applications to uncertainty quantification in subsurface flow},
% SIAM/ASA Journal on Uncertainty Quantification, 3(1):1075--1108, 2015.

% \bibitem{fairbanks}
% H.~Fairbanks, A.~Doostan, C.~Ketelsen, and G.~Iaccarino,
% \textit{A low-rank control variate for multilevel Monte Carlo 
% simulation of high-dimensional uncertain systems},
% Journal of Computational Physics, 341:121--139, 2017.

% \bibitem{nr}
% R.A.~Rossi and N.K.~Ahmed,
% \textit{The Network Data Repository with Interactive Graph Analytics 
% and Visualization},
% AAAI, 2015.

% \bibitem{leskovec2005graphs}
% J.~Leskovec, J.~Kleinberg, and C.~Faloutsos,
% \textit{Graphs over Time: Densification Laws, Shrinking Diameters and 
% Possible Explanations},
% ACM SIGKDD International Conference on Knowledge Discovery and Data 
% Mining (KDD), 2005.

% \bibitem{mcauley2012learning}
% J.~McAuley and J.~Leskovec,
% \textit{Learning to Discover Social Circles in Ego Networks},
% Advances in Neural Information Processing Systems (NIPS), 2012.


% \bibitem{crescenzi2010finding}
% P.~Crescenzi, R.~Grossi, C.~Imbrenda, L.~Lanzi, and A.~Marino,
% \textit{Finding the Diameter in Real-World Graphs: Experimentally 
% Turning a Lower Bound into an Upper Bound},
% in Algorithms -- ESA 2010, Springer Berlin Heidelberg, pp.~302--313, 
% 2010.
% \end{thebibliography}
\end{document}
